\documentclass[12pt]{jlreq}
\usepackage{amsmath, amssymb}

\newcommand{\C}{\mathbb{C}}
\newcommand{\R}{\mathbb{R}}
\newcommand{\Z}{\mathbb{Z}}
\newcommand{\N}{\mathbb{N}}
\newcommand{\F}{\mathbb{F}}
\newcommand{\Prob}{\mathbb{P}}
\newcommand{\Exp}{\mathbb{E}}
\newcommand{\dd}{\,d}
\newcommand{\Ker}{\operatorname{Ker}}
\newcommand{\Impart}{\operatorname{Im}}
\newcommand{\Hom}{\operatorname{Hom}}

\begin{document}

円周 \(S^1\) を \(\{ z \in \C \mid |z| = 1 \}\) とみなし，\(T^2 = S^1 \times S^1\) と表す．\(T^2 \times [0, 1]\) の境界の \(2\) つの成分 \(T^2 \times \{0\}\) と \(T^2 \times \{1\}\) を
\[
  (z, \zeta, 1) \sim (z, z^2 \zeta, 0), \quad (z, \zeta) \in S^1 \times S^1 = T^2
\]
に従って同一視する．\(T^2 \times [0, 1]\) からこのようにして得られる商空間を \(X\) とする．また，連続写像 \(T^2 \times [0, 1] \to T^2\)，\((z, \zeta, t) \mapsto (z, e^{2\pi\sqrt{-1}t})\) の誘導する連続写像を \(\pi : X \to T^2\) とする．

(1) \(X\) の整数係数ホモロジー群 \(H_*(X; \mathbb{Z})\) を求めよ．

(2) 連続写像 \(s : T^2 \to X\) で \(\pi \circ s = 1_{T^2}\) をみたすものが存在しないことを証明せよ．ただし \(1_{T^2}\) は \(T^2\) の恒等写像を表す．

\end{document}